\documentclass[a4paper,10pt]{report}

\def\packagepath{../../preambule}   % path du package princial
\usepackage{\packagepath/preambule}    % utilisation du fichier de configuration

\def\level{Première }              % Classe
\def\course{NSI}              % Matière
\def\eval{DS01 - Variables}

\def\visibleornot{visible}    % visible or invisible
\def\documentpath{./Sources_Latex}
\def\date{Mardi 23 Septembre 2025}
\renewcommand{\arraystretch}{1}  % Ecart dans les tableau

\def\ptsexoA{2}
\def\ptsexoB{3}
\def\ptsexoC{3}
\def\ptsexoD{3}
\def\ptsexoE{2}
\def\ptsexoF{1}
\def\ptsexoG{2}

\def\ptstotal{\ptsexoA+\ptsexoB}

\usepackage{hyperref}

\usetikzlibrary{shapes, arrows, positioning}
\usetikzlibrary{shapes.geometric, arrows, positioning, decorations.pathreplacing}

\tikzstyle{etat} = [draw, rounded corners=15pt, minimum width=2.5cm, minimum height=1.2cm, text centered, font=\bfseries]
\tikzstyle{nouveau} = [etat, fill=blue!40]
\tikzstyle{pret} = [etat, fill=green!60]
\tikzstyle{elu} = [etat, fill=yellow!60]
\tikzstyle{bloque} = [etat, fill=red!60]
\tikzstyle{termine} = [etat, fill=white]
\tikzstyle{fleche} = [thick, ->, >=stealth]

\usepackage{float} % Permet d'utiliser [H]
\usepackage[T1]{fontenc}
\begin{document}

%%%%%%%%%%%%%%%%%%%%%%%%%%%%%%%%%%%%%%%%%%%%%%%%%%ù

%\PageGardeSujetBac[Matiere = Numérique et Sciences Informatiques,
%                  Session = 2025,
%                  AffJour=false,
%                  Duree = {3 heures 30},
%                  DernierePage = \pageref{LastPage},
%                  ModeExamen=false
%                  ]
\renewcommand{\labelitemi}{\textbullet} %pour éviter les tirets dans les "itemize" qui apportent confusion avec le signe moins.
%% Début page de garde style BAC

   %%%%%%%%%%%%%%%%%%%%%%%%%%%%%%%%%%%%%%%%%%%%%%%%%%%%%%%%
%\PageGardeSujetBac[clés]

\pagestyle{DS_FP}          % Feuille de style fancy
\NomPrenomNote{}


\exods{\ptsexoA}

Indiquer si les noms de variables suivants sont pertinents ou non. Justifier votre réponse.

\medskip

\renewcommand{\arraystretch}{3.5}
\begin{tabular}{|p{3cm}|p{3cm}|p{3cm}|p{7cm}|}
        \hline
        \textbf{Nom de variable} & \textbf{Pertinent} & \textbf{Non pertinent} &\textbf{Justification} \\ \hline
    \verb|ma variable| &&X& Espace dans le nom de la variable \\ \hline
    \verb|ma_variable77|&X&& \\ \hline
    \verb|77ma_variable|&&X&Nom de variable commencant par un chiffre \\ \hline
    \verb|ma_variable@77|&&X& Présence d'un caractère spécial dans le nom de la variable\\ \hline

    \end{tabular}



\bigskip

\exods{\ptsexoB}

    \renewcommand{\arraystretch}{1}

Indiquer l'état des variables après l'exécution de chaque ligne d'instruction.

\medskip

\begin{minipage}{0.29\linewidth}

\begin{CodePythontexAlt}[TaillePolice=\large,Largeur=1\linewidth,Centre,PremLigne=1]{}
a = 3
b = 5
c = a + b
a = b - 2
b = c + 1
d = a + b + c
\end{CodePythontexAlt}

\end{minipage}\hfill
\begin{minipage}{0.66\linewidth}
    
    \medskip

    \begin{tabular}{|c|c|c|c|c|}
        \hline
        Ligne& Valeur de a & Valeur de b &Valeur de c & Valeur de d\\ \hline
        1&3&&&\\\hline
        2&&5&&\\\hline
        3&&&8&\\\hline
        4&3&&&\\\hline
        5&&9&&\\\hline
        6&&&&20\\\hline


    \end{tabular}
\end{minipage}

\bigskip

\exods{\ptsexoC}

Indiquer le type et la valeur de la variable \verb|a| après l'exécution de chacune des lignes. Si la ligne génère une erreur, indiquer "erreur".


\begin{minipage}{0.49\linewidth}

\begin{CodePythontexAlt}[TaillePolice=\large,Largeur=1\linewidth,Centre,PremLigne=1]{}
a = 3
a = a + 0.1
a = "Bonjour"
a = a + " tout le monde"
a = True
a = a + 1
a = input('Entrer un nombre : ')
a = '5'
a = a + '2'
a = int(a)
a = float(a)

\end{CodePythontexAlt}

\end{minipage}\hfill
\begin{minipage}{0.46\linewidth}
    
    \medskip

    \renewcommand{\arraystretch}{1.2}
    \begin{tabular}{|c|p{2cm}|p{5cm}|}
        \hline
        Ligne& Type de a & Valeur de a\\ \hline
        1& \verb|int|&3\\\hline
        2&float&3.1\\\hline
        3&str&"Bonjour"\\\hline
        4&str&"Bonjour tout le monde"\\\hline
        5&bool&True\\\hline
        6&ERREUR&\\\hline
        7&str&\\\hline
        8&str&'5'\\\hline
        9&str&'52'\\\hline
        10&int&52\\\hline
        11&float&52.0\\\hline


    \end{tabular}
\end{minipage}


\pagebreak
\pagestyle{DS_LP}
\exods{\ptsexoD}

On considère le code suivant:

\begin{CodePythontexAlt}[TaillePolice=\large,Largeur=1\linewidth,Centre,PremLigne=1]{}
mot = 'Numerique'
\end{CodePythontexAlt}


Indiquer ce qui est affiché par les instructions suivantes. Si cela génère une erreur, indiquer "erreur".

\renewcommand{\arraystretch}{1.2}
    \begin{tabular}{|p{3cm}|p{4cm}|}
        \hline
         Instruction & Affichage\\ \hline
         \verb|mot[0]|&'N'\\\hline
        \verb|mot[2]| &'m'\\\hline
        \verb|mot[4]| &'r'\\\hline
        \verb|mot[8]|&'e'\\\hline
        \verb|len(mot)|&9\\\hline
        \verb|mot + 'nsi'| &'Numeriquensi'\\\hline
        
        
    \end{tabular}\hfill
\renewcommand{\arraystretch}{1.2}
    \begin{tabular}{|p{3cm}|p{4cm}|}
        \hline
         Instruction & Affichage\\ \hline
        \verb|'nsi' + mot|&'nsiNumerique'\\\hline
        \verb|'u' in mot|&True\\\hline
        \verb|'n' in mot|&False\\\hline
        \verb|int(mot)|&ERREUR\\\hline
        \verb|mot * 2| &'NumeriqueNumerique'\\\hline
        \verb|mot[3] * 2| &'ee'\\\hline
        
        


    \end{tabular}

\bigskip

\exods{\ptsexoE}

On considère le code suivant:

\begin{CodePythontexAlt}[TaillePolice=\large,Largeur=1\linewidth,Centre,PremLigne=1]{}
a = True
b = False
\end{CodePythontexAlt}


Indiquer ce qui est affiché par les instructions suivantes. Si cela génère une erreur, indiquer "erreur".

\renewcommand{\arraystretch}{1.2}
    \begin{tabular}{|p{4cm}|p{3cm}|}
        \hline
         Instruction & Affichage\\ \hline
         \verb|a and b|&False\\\hline
        \verb|a or b| &True\\\hline
        \verb|not a| &False\\\hline
        \verb|a and (not b)|&True\\\hline
        
        
    \end{tabular}\hfill
\renewcommand{\arraystretch}{1.2}
    \begin{tabular}{|p{4cm}|p{3cm}|}
        \hline
         Instruction & Affichage\\ \hline
        \verb|a or (not b)| &True\\\hline
        \verb|a and (a or b)|&True\\\hline
        \verb|b or (a and b)|&False\\\hline
        \verb|a and (not(a and b))|&True\\\hline
    \end{tabular}
\bigskip


\exods{\ptsexoF}

On considère le code suivant:

\begin{CodePythontexAlt}[TaillePolice=\large,Largeur=1\linewidth,Centre,PremLigne=1]{}
a = 3
b = 4
\end{CodePythontexAlt}


Indiquer ce qui est affiché par les instructions suivantes. Si cela génère une erreur, indiquer "erreur".

\renewcommand{\arraystretch}{1.2}
    \begin{tabular}{|p{4cm}|p{3cm}|}
        \hline
         Instruction & Affichage\\ \hline
         \verb|a == 3|&True\\\hline
        \verb|a != 3| &False\\\hline
        \verb|a > 3| &False\\\hline
        
        
        
    \end{tabular}\hfill
\renewcommand{\arraystretch}{1.2}
    \begin{tabular}{|c|p{4cm}|p{3cm}|}
        \hline
         Instruction & Affichage\\ \hline
        \verb|type(a) == str|&False\\\hline
        \verb|a == 4 and a/0 == 1|&False\\\hline
        \verb|a == 3 or a/0 == 1| &True\\\hline
         \end{tabular}
\bigskip


\exods{\ptsexoG}

Proposer un code Python permettant:

\begin{enumerate}
    \item de demander à un utilisateur son année de naissance.
    \item de lui demander son prénom
    \item de lui afficher le message suivant: "Bonjour \texttt{<prenom>}, tu as \texttt{<age>} ans." en remplaçant \verb|<age>| et \verb|<prenom>| par l'âge et le prénom de l'utilisateur. Par exemple:  \quad \verb|Bonjour Thomas, tu as 17 ans.|
    \end{enumerate}
    
\begin{CodePythontexAlt}[TaillePolice=\large,Largeur=1\linewidth,Centre,PremLigne=1]{}
anne = input('Entrer votre année de naissance : ')
prenom = input('Entrer votre prénom : ')
age = 2025 - int(anne)
print('Bonjour ' + prenom + ', tu as ' + str(age) + ' ans.')
\end{CodePythontexAlt}
\end{document}








